\section{Discussion}
\label{sec:discussion}

\subsection{Two distinct small-message bottlenecks}

Section~\ref{sec:results} shows the same latency-bound floor at small message sizes across every algorithm variant, and the cost model of Equation~\ref{eq:hockney} explains why in a way that separates two mechanisms which are easy to conflate into a single vague "small messages are slower."

\textbf{Pipeline fill latency} is structural. With only $N$ chunks and no further sub-chunking, every one of the $2(N-1)$ steps pays the full per-message latency $\alpha$ regardless of how small $M$ is; that fixed cost does not amortize over a smaller payload the way the bandwidth term does. As $M \to 0$, $\text{busbw} = M/T$ collapses toward zero simply because $T$ floors out at $2(N-1)\alpha$ while the numerator keeps shrinking. This is the flat region every panel of Figure~\ref{fig:busbw-vs-size} shows, and Figure~\ref{fig:latency-decomposition} confirms it directly: at the smallest swept size the modeled bandwidth term is negligible next to the modeled latency term at every $N$. The magnitude is concrete rather than abstract. With the custom ring's fitted $\alpha = 0.264~\mu$s, the predicted floor at $N=16$ is $2(16-1) \times 0.264 \approx 7.9~\mu$s, which is the right order for the 11.6~$\mu$s actually measured at 8 bytes, the remainder being the synchronization cost discussed next.

\textbf{Rank synchronization overhead} is a different, and non-structural, cost sitting on top of that floor. Step $s+1$ on rank $r$ cannot start until rank $r$'s own step $s$ has completed \emph{and} its neighbor has actually sent, so the operation runs at the speed of its slowest participant, which is exactly what the \texttt{MPI\_MAX} reduction of Section~\ref{sec:experimental-setup} measures. Stragglers and scheduling jitter therefore inflate the measured time on top of whatever the hardware's genuine point-to-point latency is. Figure~\ref{fig:latency-decomposition} is what makes this measurable rather than asserted, by putting the measured 8-byte time next to an ideal floor of $2(N-1)$ raw point-to-point latencies (0.1585~$\mu$s each, measured on the same hardware in the same run).

The answer is more interesting than a single constant. Averaged over the whole sweep, the fit puts the ring's per-step $\alpha$ at 0.264~$\mu$s against that 0.1585~$\mu$s bare latency, a gap of roughly 0.105~$\mu$s per step. But the per-$N$ view shows that average is hiding a threshold. Up to about $N=14$ the measured time sits essentially \emph{on} the ideal floor: the implementation's per-step cost really is close to a bare point-to-point send, and synchronization overhead is small. At $N=15$ and $N=16$ it departs sharply, reaching 11.55~$\mu$s against a 4.7~$\mu$s floor. So the honest statement is not "the ring pays a fixed 0.105~$\mu$s of synchronization per step" but "the ring pays almost no synchronization overhead until rank placement turns against it, and then it pays a great deal."

The most likely mechanism is this CPU's hybrid topology: 8 performance cores with two hardware threads each, plus 12 slower efficiency cores. Nothing is oversubscribed at $N=16$ on a 28-thread machine, so this is not classical core contention. It is that the \emph{slowest} rank sets the pace, and past a certain count some rank inevitably lands on a hyperthread sibling or an efficiency core and becomes a straggler that all $2(N-1)$ steps then wait on. This is worth flagging as a property of the measurement platform rather than of ring-allreduce: on a homogeneous cluster, or with explicit rank binding to performance cores only, the effect would be expected to weaken or disappear, and confirming that is a natural follow-up.

\subsection{Why a correct, bandwidth-optimal ring still loses at small messages}

Section~\ref{sec:step-count}'s step-count table is the other half of the story, and the three-variant comparison of Section~\ref{sec:head-to-head} turns it from an argument into a measurement.

Ring-allreduce is bandwidth-optimal, its total data movement is asymptotically independent of $N$, but it is not \emph{step}-optimal: it pays $2(N-1)$ sequential round trips where an algorithm built on recursive halving and doubling, such as the one Rabenseifner describes and \citet{thakur2005optimization} analyze, pays only $O(\log N)$. At $N=16$ that is 30 steps against 8, and every one of those extra steps costs at least one $\alpha$. Open~MPI's \texttt{coll/tuned} component exposes exactly those alternatives as separately selectable algorithms (Section~\ref{sec:mca-forcing} lists \texttt{recursive\_doubling} and \texttt{rabenseifner} in its enumerator), and its default selection is free to pick them for whatever size and process count it judges best.

The tempting but wrong conclusion from "my ring is five times slower than \texttt{MPI\_Allreduce} at 8 bytes" is that the implementation is poor. The forced-ring variant refutes that directly. When Open~MPI is made to run \emph{the same algorithm}, its own ring costs 11.97~$\mu$s at $N=16$ and 64~bytes, against this project's 11.08~$\mu$s: the two rings agree to within about 8\%, and both are roughly five times slower than the default selection. The gap to the default is therefore not an implementation gap at all. It is the price of the algorithm, paid identically by a mature vendor implementation and a from-scratch one, and it is precisely the $2(N-1)$ versus $O(\log N)$ step count that Table~\ref{tab:step-count} predicts. Benchmarking only against the default would have produced a number that looks like a verdict on code quality and is nothing of the sort.

Where the two rings do differ is instructive in the opposite direction. Table~\ref{tab:fits}'s fitted constants say the custom ring has the worse per-message latency ($\alpha = 0.264~\mu$s against the vendor ring's $0.113~\mu$s) but the better inverse bandwidth ($1/\beta = 8.06$~GB/s against 6.61~GB/s), and the head-to-head confirms it: the vendor ring wins the smallest sizes, the custom ring wins from 64~KiB upward, completing 16~MiB at $N=16$ in 19.2~ms against 24.6~ms. A plausible reading is that Open~MPI's ring is more careful about per-message setup while this implementation's straightforward, contiguous, non-blocking chunk exchange moves bulk bytes with less overhead per byte. That is a real and defensible outcome for a from-scratch implementation: competitive with, and at large messages faster than, the vendor's own ring, while paying more for each of its steps.

\subsection{Limitations}

Every measurement this report analyzes is single-node, gathered on the one workstation of Section~\ref{sec:experimental-setup} under Open~MPI's shared-memory transport. Two consequences follow and should be stated plainly. First, the cache-resident busbw peak, and the efficiency band that rises above 1 before falling back, are artifacts of shared memory: they would not appear over a real network fabric, where per-byte cost is set by the link rather than by whether the working set fits in L3. This is also why the linear cost model fits only loosely here ($R^2$ between 0.28 and 0.33), and a network-bound sweep would be expected to track it far better. Second, the fitted $\alpha$ folds in shared-memory handshake and scheduling costs that a network's own latency would replace rather than merely add to, so the 0.105~$\mu$s synchronization overhead quantified above is specific to this transport, even though the mechanism producing it is not.

A real multi-node sweep is therefore the single most valuable extension, and \texttt{scripts/run\_slurm\_sweep.sbatch} exists so this project's methodology extends to that setting without modification. Nothing about the algorithm, the cost model, or the analysis pipeline is single-node-specific; only the dataset gathered so far is.

The comparison is also CPU-only. GPU-to-GPU collectives, NCCL's actual production setting, add device-to-host transfer costs, NVLink or PCIe topology, and GPUDirect-style optimizations that a CPU MPI benchmark cannot speak to; that comparison is out of scope here and left to Section~\ref{sec:conclusion}'s future work rather than attempted with a mismatched proxy.

\subsection{What a segmented ring would change}

The pipeline fill latency mechanism above suggests a concrete next optimization, not attempted in this implementation: sub-chunking each of the $N$ pieces further and pipelining transmission of those sub-chunks across steps, so that a step's communication for one sub-chunk can overlap with communication for another instead of every step waiting on one full-size round trip. That directly attacks the $2(N-1)\alpha$ term, which this report's own data identifies as the ring's actual handicap: it is what makes the ring lose to the default selection at small messages, and it is the fixed cost a lower-step algorithm avoids. The bandwidth term is where this ring already wins, beating even the vendor's own ring above 64~KiB, so latency, not byte-moving, is the part still worth optimizing. That Open~MPI ships a \texttt{segmented\_ring} algorithm as a separate enumerator value (Section~\ref{sec:mca-forcing}) is a strong hint that this is the standard next move, and Section~\ref{sec:conclusion} lists it as the highest-value follow-up specifically because this report's data, not intuition, is what points at it.
